\documentclass[10pt]{scrartcl}
\usepackage[a4paper, total={5.5in, 8in}]{geometry}
\usepackage[usenames,dvipsnames,svgnames]{xcolor}
\usepackage[shortlabels]{enumitem}
\usepackage[framemethod=TikZ]{mdframed}
\usepackage{amsmath,amssymb,amsthm}
\usepackage[colorlinks]{hyperref}

\hypersetup{
	colorlinks=true,
	linkcolor=blue,
	filecolor=blue,      
	urlcolor=blue,
	pdftitle={MAT 320 Notes}
}

\usepackage{microtype}
\usepackage{mathtools}
\usepackage[headsepline]{scrlayer-scrpage}
\usepackage{thmtools}
\usepackage[textsize=small]{todonotes}

\addtolength{\textheight}{3.14cm}
\providecommand{\ol}{\overline}
\providecommand{\eps}{\varepsilon}
\providecommand{\half}{\frac{1}{2}}
\providecommand{\dang}{\measuredangle} %% Directed angle
\providecommand{\CC}{\mathbb C}
\providecommand{\FF}{\mathbb F}
\providecommand{\NN}{\mathbb N}
\providecommand{\QQ}{\mathbb Q}
\providecommand{\RR}{\mathbb R}
\providecommand{\ZZ}{\mathbb Z}
\providecommand{\dg}{^\circ}
\providecommand{\ind}[1]{\mathbf 1_{#1}}
\providecommand{\ii}{\item}
\providecommand{\alert}[1]{{\sffamily\textbf{\textcolor{magenta}{#1}}}}
\providecommand{\opname}{\operatorname}
\providecommand{\li}{\opname{li}}
\providecommand{\ls}{\opname{ls}}
\providecommand{\inv}{^{-1}}
\providecommand{\setdiff}{\smallsetminus}
\providecommand{\mb}[1]{\mathbf{#1}}
\providecommand{\abs}[1]{\left|{#1}\right|}
\providecommand{\dif}{\;d}

\reversemarginpar
\setlength{\marginparwidth}{3cm}

\mdfdefinestyle{mdbluebox}{roundcorner=10pt,innerbottommargin=9pt,
	linecolor=blue,backgroundcolor=TealBlue!5,}
\declaretheoremstyle[headfont=\sffamily\bfseries\color{MidnightBlue},
mdframed={style=mdbluebox},]{thmbluebox}

\mdfdefinestyle{mdredbox}{frametitlefont=\bfseries,innerbottommargin=8pt,
	nobreak=true,backgroundcolor=Salmon!5,linecolor=RawSienna,}
\declaretheoremstyle[headfont=\bfseries\color{RawSienna},
mdframed={style=mdredbox},headpunct={\\[3pt]},postheadspace=0pt,]{thmredbox}

\mdfdefinestyle{mdgreenbox}{linecolor=ForestGreen,backgroundcolor=ForestGreen!5,
	linewidth=2pt,rightline=false,leftline=true,topline=false,bottomline=false,}
\declaretheoremstyle[headfont=\bfseries\sffamily\color{ForestGreen!70!black},
mdframed={style=mdgreenbox},headpunct={ --- },]{thmgreenbox}

\mdfdefinestyle{mdorangebox}{linecolor=RedOrange,backgroundcolor=RedOrange!5,
	linewidth=2pt,rightline=false,leftline=true,topline=false,bottomline=false,}
\declaretheoremstyle[headfont=\bfseries\sffamily\color{RedOrange!70!black},
mdframed={style=mdorangebox},headpunct={ --- },]{thmorangebox}

\mdfdefinestyle{mdblackbox}{linecolor=black,backgroundcolor=RedViolet!5!gray!5,
	linewidth=3pt,rightline=false,leftline=true,topline=false,bottomline=false,}
\declaretheoremstyle[mdframed={style=mdblackbox}]{thmblackbox}


\declaretheorem[style=thmredbox,name=Problem,numberwithin=section]{problem}

\declaretheorem[style=thmredbox,name=Example,sibling=problem]{example}
\declaretheorem[style=thmredbox,name=$\bigstar$ Problem,sibling=problem]{niceprob}

\declaretheorem[style=thmbluebox,name=Theorem,sibling=problem]{theorem}
\declaretheorem[style=thmbluebox,name=Corollary,sibling=theorem]{corollary}
\declaretheorem[style=thmbluebox,name=Proposition,sibling=theorem]{prop}
\declaretheorem[style=thmbluebox,name=Lemma,sibling=theorem]{lemma}
\declaretheorem[style=thmbluebox,name=Theorem,numbered=no]{theorem*}
\declaretheorem[style=thmbluebox,name=Lemma,numbered=no]{lemma*}
\declaretheorem[style=thmgreenbox,name=Claim,numberwithin=problem]{claim}
\declaretheorem[style=thmgreenbox,name=Claim,numbered=no]{claim*}
\declaretheorem[style=thmblackbox,name=Remark,sibling=theorem]{remark}
\declaretheorem[style=thmblackbox,name=Remark,numbered=no]{remark*}
\declaretheorem[style=thmgreenbox,name=Definition,sibling=theorem]{definition}
\declaretheorem[style=thmgreenbox,name=Definition,numbered=no]{definition*}
\declaretheorem[style=thmorangebox,name=Warning,sibling=theorem]{warning}
\declaretheorem[style=thmorangebox,name=Warning,numbered=no]{warning*}
\declaretheorem[style=thmorangebox,name=Note,numbered=no]{note}
\declaretheorem[style=thmblackbox,name=Exercise,sibling=example]{exercise}
\declaretheorem[style=thmblackbox,name=Exercise,numbered=no]{exercise*}

\newenvironment{soln}{\begin{proof}[Solution]}{\end{proof}}
\newenvironment{walkthrough}{\noindent \textbf{Walkthrough.}}{\medskip}
\newlist{walk}{enumerate}{3}
\setlist[walk]{label=\bfseries (\alph*)}

\providecommand{\pow}{\operatorname{Pow}}
\providecommand{\vocab}[1]{{\textbf{\textcolor{ForestGreen}{#1}}}}
\providecommand{\lcm}{\operatorname{lcm}}
\providecommand{\ord}{\operatorname{ord}}

\usepackage{asymptote}
\begin{asydef}
	defaultpen(fontsize(10pt));
	unitsize(1cm);
	size(8cm); // set a reasonable default
	usepackage("amsmath");
	usepackage("amssymb");
	settings.tex="pdflatex";
	settings.outformat="pdf";
	// Replacement for olympiad+cse5 which is not standard
	import geometry;
	// recalibrate fill and filldraw for conics
	void filldraw(picture pic = currentpicture, conic g, pen fillpen=defaultpen, pen drawpen=defaultpen)
	{ filldraw(pic, (path) g, fillpen, drawpen); }
	void fill(picture pic = currentpicture, conic g, pen p=defaultpen)
	{ filldraw(pic, (path) g, p); }
	// some geometry
	pair foot(pair P, pair A, pair B) { return foot(triangle(A,B,P).VC); }
	pair orthocenter(pair A, pair B, pair C) { return orthocentercenter(A,B,C); }
	pair centroid(pair A, pair B, pair C) { return (A+B+C)/3; }
	// cse5 abbreviations
	path CP(pair P, pair A) { return circle(P, abs(A-P)); }
	path CR(pair P, real r) { return circle(P, r); }
	pair IP(path p, path q) { return intersectionpoints(p,q)[0]; }
	pair OP(path p, path q) { return intersectionpoints(p,q)[1]; }
	path Line(pair A, pair B, real a=0.6, real b=a) { return (a*(A-B)+A)--(b*(B-A)+B); }
	// cse5 more useful functions
	picture CC() {
		picture p=rotate(0)*currentpicture;
		currentpicture.erase();
		return p;
	}
	pair MP(Label s, pair A, pair B = plain.S, pen p = defaultpen) {
		Label L = s;
		L.s = "$"+s.s+"$";
		label(L, A, B, p);
		return A;
	}
	pair Drawing(Label s = "", pair A, pair B = plain.S, pen p = defaultpen) {
		dot(MP(s, A, B, p), p);
		return A;
	}
	path Drawing(path g, pen p = defaultpen, arrowbar ar = None) {
		draw(g, p, ar);
		return g;
	}
\end{asydef}

\ihead{\footnotesize\textbf{MAT 320 Fall 2024 Notes}}
\ohead{\footnotesize \today}

\title{\vspace{-2.0cm} MAT 320 Fall 2024}
\author{\large Aditya Pahuja}
\date{\large \today}
\begin{document}
\maketitle
\tableofcontents
\pagebreak

These are notes for MAT 320 only (i.e. starting week 7).

\setcounter{section}{6}
\section{Week 7: 10/08 and 10/10}
The following two theorems are some of the most important
in analysis over $\RR$.
\begin{theorem}[Weierstrass theorem]\label{wk7:thm:weierstrass}
	If $f\colon\RR\to\RR$ is a continuous function on the interval $[a, b]$,
	then it has a maximum and a minimum.
	More precisely, the set $\{f(x) : x\in[a,b]\}$ has a minimum and a maximum.
\end{theorem}

\begin{theorem}[Intermediate value theorem]\label{wk7:thm:ivt}
	If $f\colon\RR\to\RR$ is continuous on an interval $I$
	(whether open, half-open, or closed)
	and $a < b$ are two elements of $I$, then,
	if $y$ is between $f(a)$ and $f(b)$
	(so either $f(a) < y < f(b)$ or $f(b) < y < f(a)$),
	there exists an $x\in(a,b)$ such that $f(x) = y$.
\end{theorem}

\begin{exercise}
	Show that the intermediate value theorem equivalently says that
	$f(I)$ is also an interval.
\end{exercise}

Some toy applications of these theorems:
\begin{theorem*}
	Let $y\ge 0$ be real and $m$ an integer.
	Then, there exists a real $x\ge 0$ such that $x^m = y$.
\end{theorem*}

Note that there is no uniqueness assertion here!
\begin{proof}
	Let $f\colon \RR\to\RR$ be the function $f(x) = x^m$.
	We check three cases.
	\begin{itemize}
		\ii If $y = 0$ or $y = 1$, then the result is trivial, as
		we take $x = 0$, $x = 1$ respectively.
		\ii If $0 < y < 1$, then $f(0) < y < f(1)$ implies that
		there exists an $x\in(0,1)$ such that $f(x) = y$,
		by the intermediate value theorem.
		\ii If $1 < y$, then $f(1) < y < f(y)$ implies that
		there exists an $x\in(1,y)$ such that $f(x) = y$,
		again by the intermediate value theorem.
	\end{itemize}
	These cases encompass all possible values of $y$, so we are done.
\end{proof}
Note in fact that we got an estimate on how large $x$ must be
for a given $y$, too! (A bad estimate, but still.)

Some more problems for the reader:
\begin{problem}
	Show that if $f\colon[0,1]\to[0,1]$ is continuous
	then $f$ has a fixed point;
	that is, there exists an $x\in[0,1]$ such that $f(x) = x$.
\end{problem}

\begin{problem}
	On Saturday, a mountain climber climbs from the base of a mountain
	to its peak between 7 am and 10 pm.
	On Sunday, the mountain climber climbs down from the peak to the base,
	again between 7 am and 10 pm.
	Prove that at some point on Sunday, the mountain climber
	is at the same altitude that he was at exactly 24 hours ago.
\end{problem}

Rather than just proving these theorems for real numbers,
we will depart into the world of \emph{metric spaces}
to help us understand them in a much deeper context.

\subsection{Metric spaces}
Fundamentally, there is really only one thing that matters
in the context of the basic ideas of convergence
(which is what basically everything so far has been based on):
the notion of a distance. In light of this,
we momentarily strip away all the extra features of $\RR$
to get the following.

\begin{definition}[Metric space]\label{wk7:def:metricspace}
	A \vocab{metric space} is a set $X$
	equipped with a \vocab{metric} (``distance function'')
	$d\colon X\times X\to \RR$ such that, for all $x,\,y,\,z\in X$,
	\begin{itemize}
		\ii $d(x,y)\ge 0$, with equality if and only if $x = y$;
		\ii $d(x, y) = d(y, x)$; and
		\ii $d(x, y) \le d(x, z) + d(z, x)$.
	\end{itemize}
	We refer to metric spaces as a pair $(X, d)$,
	or, if $d$ is clear from context, then we will just say $X$.
	If we just say ``$X$ is a metric space,''
	then the corresponding metric will be called $d$.
	We also often refer to the elements of $X$ as ``points''
	to align with visual intuition.
\end{definition}

\begin{example}
	Some example metric spaces:
	\begin{itemize}
		\ii Any subset of the real numbers is a metric space
		with metric $d(x, y) = |x - y|$.
		\ii In general, $n$-dimensional space is a metric space
		with the Euclidean metric
		\[d(\mb x, \mb y) = \sqrt{(x_1-y_1)^2 + \cdots + (x_n - y_n)^2}\]
		or with the taxicab metric
		\[d(\mb x, \mb y) = |x_1 - y_1| + \cdots + |x_n - y_n|\]
		or with the Chebyshev metric
		\[d(\mb x, \mb y) = \max\{|x_1 - y_1|, \dots, |x_n - y_n|\}.\]
		\ii A (connected) graph can be turned into a metric space
		with distance given by the length of the shortest path
		between two points.
		\ii We also have the trivial metric
		\[d(x, y) = \begin{cases}
			0 & \text{if }x=y\\
			1 & \text{if }x\ne y.
		\end{cases}\]
		\ii The rationals can also be made a metric space
		with distance $d(a,b) = p^{-\nu_p(a-b)}$;
		this is called the $p$-adic metric.
		Here $\nu_p(n)$ is the largest power of $p$ dividing $n$
		if $n$ is a nonzero integer, $\infty$ if $n = 0$,
		and we extend to rationals by decreeing that
		$\nu_p(a/b) = \nu_p(a) - \nu_p(b)$ for all rationals;
		note that we define $p^{-\infty} = 0$.
		Actually, the rationals with this metric form an ultrametric space,
		meaning the stronger version of the triangle inequality
		\[d(x,y) \le \max\{d(x, z), d(y, z)\}\]
		always holds (this is not too hard to prove
		if you know some number theory).
	\end{itemize}
\end{example}

Our job now is to carry over our basic definitions
from $\RR$ to general metric spaces.
Recalling that $|x - y|$ just means ``the distance between $x$ and $y$''
in $\RR$, this is just a matter of copying the definitions and
changing absolute values into distances. For example:

\begin{definition}[Convergence]\label{wk7:def:convergence}
	Let $(X, d)$ be a metric space
	and let $(x_n)$ be a sequence of points in $X$.
	Then, we say that $(x_n)$ converges to a point $\ell\in X$ if
	\[\text{for all }\eps > 0,\quad d(x_n,\ell) < \eps\text{ eventually.}\]
\end{definition}

Similarly, we can define what it means for a function to be continuous;
note that since we compare distances both in the domain (``$|x - c| < \delta$'')
and in the codomain (``$|f(x) - f(x)| < \eps$''),
both the domain and codomain need to be metric spaces
(but not necessarily the same one!).

\begin{definition}[Continuity]\label{wk7:def:continuity}
	A function $f\colon X\to Y$ between metric spaces $(X, d_X)$ and $(Y, d_Y)$
	is continuous at a point $c\in X$ if
	\[\text{for all }\eps > 0,\text{ there exists }\delta > 0
	\text{ such that }d_X(x, c) < \eps \implies d_Y(f(x), f(c)) < \eps.\]
\end{definition}

\subsection{Some topology}

At this point, you might have made an observation:
the idiom $d(a, x) < \eps$ keeps popping up,
since we frequently need to refer to the idea of
``$x$ is less than $\eps$ away from $a$.''
To make notation less clunky, we define:
\begin{definition}\label{wk7:def:openball}
	In a metric space $(X, d)$, the \vocab{open ball}
	centered at $x\in X$ with radius $\eps > 0$ is the set of points
	\[B_X(x, \eps) := \{y\in X : d(x, y) < \eps\}.\]
	The \vocab{open unit ball} is then just $B_X(x, 1)$.
	We sometimes omit the subscript (simply writing $B(x,\eps)$)
	if there is no risk of confusion about the ambient metric space.
\end{definition}

The adjective ``open'' here refers to the idea that this set
doesn't have its natural ``boundary'' of points
whose distance from $x$ is exactly $\eps$ ---
the open balls in $\RR^n$ should give you a clearer idea of what this means.
We will see a slightly more specific example of what this means
in a moment.

\begin{exercise}
	Pick your three favorite metric spaces.
	What do open balls look like in those spaces?
\end{exercise}

\begin{exercise}
	Rewrite Definitions \ref{wk7:def:continuity} and \ref{wk7:def:convergence}
	with the notation of open balls.
\end{exercise}

The kicker is that we very much care about sets $U\subseteq X$
such that, for every $x\in U$, some open ball $B(x,\eps)$ with $\eps > 0$
is contained in $U$, since these are contexts in which
we can freely talk about convergence, continuity, and so on
(we only require existence of an $\eps$
because we only ever care about ``arbitrarily small $\eps$''
In fact, such $U$ are so important that we give them a name:
\begin{definition}[Open and closed sets]\label{wk7:def:openclosed}
	Let $(X, d)$ be a metric space as usual.
	A subset $U\subseteq X$ is said to be \vocab{open}
	(in $X$, with respect to the metric $d$ ---
	both of these descriptors are dropped if we have enough context)
	if, for every $x\in U$, there exists an $\eps$ such that $B(x,\eps)$
	is a subset of $U$.
	
	The set $E\subseteq X$ is \vocab{closed} if its complement
	$X\setdiff E$ is open.
\end{definition}

\begin{exercise}
	Show that every open ball is open in its respective metric space,
	this earning open balls their name.
\end{exercise}

\begin{exercise}
	Let $(X, d)$ be a metric space.
	Prove the following facts about open sets:
	\begin{enumerate}[(i)]
		\ii $\varnothing$ and $X$ are open in $X$.
		\ii The union of any (possibly infinite) collection
		of open sets is also an open set.
		\ii The intersection of finitely many open sets is also an open set.
		Show that finiteness is a necessary assumption here
		by finding an infinite collection of open sets in $\RR$
		whose intersection is not open in $\RR$.
	\end{enumerate}
	What are the analogous statements for closed sets?
	Write them down and prove them, too.
\end{exercise}

Closed sets are in fact aptly named too:
\begin{theorem}[Alternate definition of closed set]
	\label{wk7:def:closed-iff-contains-limitpts}
	Let $E\subseteq X$.
	A point $x\in X$ is a \vocab{limit point} of $E$
	if there exists a sequence $(x_n)$ of points in $E$
	whose limit is $x$.
	Then, $E$ is closed if and only if
	it contains all of its limit points. 
\end{theorem}

\begin{proof}
	\todo{proof}
\end{proof}

You can think of this as ``closed sets do contain their boundaries.''
\begin{exercise}
	Show that finite sets are closed in every metric space.
	Then, find a metric space in which all finite sets are open,
	and find a metric space in which no finite sets are open.
\end{exercise}

The following theorem justifies the above definitions
by allowing us to think about continuity much more globally
rather than pointwise
(in fact, this will be the exact definition of continuity
in general topological spaces):
\begin{theorem}[Continuity, topologically]
	A function $f\colon X\to Y$ between metric spaces
	is continuous if and only if, for each open $U\subseteq Y$,
	the set $f\inv(U)\subseteq X$ is also open.
	This theorem also holds if the word ``open'' is replaced with ``closed.''
\end{theorem}

The proof of this is given in Week 8,
and a formal writeup was required for homework 7.
You should try to draw a picture and prove the theorem yourself
before reading on; being able to prove it yourself
means you can be confident in your understanding of the preceding content.

\section{Week 8: 10/17}
From the end of last week's lecture:
\begin{theorem}[Continuity, topologically]
	\label{wk8:thm:topocont}
	Let $f\colon X\to Y$ be a function between two metric spaces.
	Prove that $f$ is continuous if and only if,
	for every open subset $U\subseteq Y$, the pre-image
	$f\inv(U)$ is open in $X$.
\end{theorem}

\begin{proof}
	Here's a picture:
	\begin{center}
		\begin{asy}
		unitsize(1.15cm);
		draw((1.2,0)..(0,2)..(-1.2,0)..(0,-2)..cycle);
		draw((6.2,0)..(5,2)..(3.8,0)..(5,-2)..cycle);
		pair x = (0.3,0.5);
		filldraw((0.3,0.08)..(0.6,0.2)..(0.72,0.37)..(0.3,1.5)..(-0.6,1.3)..(-0.4,0.2)..cycle,green+opacity(0.2),dashed);
		filldraw(circle(x,0.35),red+opacity(0.3),dashed);
		dot("$x$", x, dir(90));
		pair fx = (4.8, -.3);
		filldraw(circle(fx,0.7),green+opacity(0.2),dashed);
		filldraw((5.1,-.34)..(4.8,-0.1)..(4.5,0.2)..(4.4,-0.2)..(5.2,-0.7)..cycle,red+opacity(0.3),dashed);
		dot("$f(x)$", fx, dir(90),fontsize(9));
		draw((0.7,0.5)..(2.5,0.3)..(4.3,-0.2),EndArrow(5));
		
		label("$B_X(x,\delta)$",x+(0,0.3),dir(100),fontsize(7.5));
		label("$B_Y(f(x),\varepsilon)$",fx-(0,0.7),dir(-80),fontsize(7.5));
		label("$f^{-1}(U)$",(0,2),dir(90),fontsize(11));
		label("$U$",(5,2),dir(90),fontsize(11));
		label("$f$",(2.5,0.3),dir(90));
		label("$f^{-1}$",(2.5,0.15),dir(-90));
		
		draw((0.8,0.35)..(2.5,0.15)..(4.05,-0.28),BeginArrow(5));
		label("$f^{-1}(B_Y(f(x),\varepsilon))$",(0,0.1),dir(-90),fontsize(7.5));
		\end{asy}
	\end{center}
	
	First, suppose $f$ is continuous.
	For an open $U\subseteq Y$, let $x\in f\inv(U)$
	and let $B_Y(f(x),\eps)$ be contained in $U$.
	Then, there exists a $\delta > 0$ so that
	$f(B_X(x,\delta))\subseteq B_Y(f(x),\eps)$ by definition of continuity.
	This means that $B_X(x,\delta)\subseteq f\inv(U)$, i.e.
	for each $x\in f\inv(U)$, we can find an open ball centered at $x$
	contained in $U$, which is what it means for $f\inv(U)$ to be open.
	
	Conversely, suppose $f\inv(U)$ is open in $X$ for all open $U\subseteq Y$.
	Then, let $x$ be in $Y$ and consider the set $f\inv(B_Y(f(x),\eps))$.
	The ball is open in $Y$, so this pre-image is open in $X$
	and contains $f\inv(x)$. This means we can find an open ball
	$B_X(x,\delta)\subseteq f\inv(U)$ due to openness,
	but then the existence of $\delta$ is exactly what is required
	to prove $f$ is continuous, so we are done.
\end{proof}

\subsection{Subspaces of metric spaces}
Let $(X, d)$ be a metric space.
There are a couple of natural ways to generate more metric spaces
given a few metric spaces: for example, homework 6 problem 6
showed that the product of two metric spaces is a metric space
in a natural way.
We can also create a metric space by taking a subset
and just restricting the metric to the subset; that is,
if $E\subseteq X$, then we can define a metric on $E$ by
\[d_{E} = d|_{E\times E}.\]

For example, let $E$ be the unit circle $\{(x, y) : x^2 + y^2 = 1\}$.
Then, we can define the distance between two points by\dots
well, their distance (in the standard Euclidean sense).

\begin{center}
	\begin{asy}
		draw((-3,0)--(3,0));
		draw((0,-3)--(0,3));
		draw(circle((0,0),2));
		pair A = 2*dir(30);
		pair B = 2*dir(140);
		draw(A--B, dashed);
		dot(A);
		dot(B);
	\end{asy}
\end{center}

Or, for a simpler toy example, we can take $E = (0,1]\subseteq\RR$.
Then, \emph{within $E$}, an open ball $B_E(z, \eps)$ centered at $z$
with radius $\eps$ would be
\[B_E(z, \eps) = \{x\in E: d_E(x, z) < \eps\} =
\{x\in X : d(x, z) < \eps\} \cap E = B_X(z,\eps) \cap E.\]
That is, to look at an open ball in $E$,
we take the corresponding open ball in $X$ and chop off
the parts outside $E$ (and this generalizes to all open sets).
Crucially, this means that \alert{openness is unique to each metric space}:
a set can be open in $E$ while not being open in $X$.
In our example with $X = \RR$, $E = (0, 1]$, $B_E(1, \eps) = (1-\eps,1]$
is open in $E$, but not open in $\RR$
(because any open ball in $\RR$ centered at $1$ would contain
numbers bigger than $1$).

\begin{exercise}
	What do open sets look like on the unit circle,
	when treated as a subset of $\RR^2$ with the usual distance metric?
\end{exercise}

\subsection{Connectedness}

\begin{definition}
	\label{wk8:def:connected}
	A metric space $(X, d)$ is \vocab{connected} if
	there are no nontrivial sets that are both open and closed (clopen).
	That is, if $A\subset X$ is clopen, then $A = X$ or $A = \varnothing$.
	We say that $E\subseteq X$ is connected if
	$E$ is connected with respect to the metric inherited from $X$.
	
	If the space is not connected then we say it is \vocab{disconnected}.
\end{definition}

Perhaps the more natural way to think of this is that
$X$ is connected if it cannot be expressed as
the union of two nonempty disjoint open sets.
This is equivalent to the definition because
if $X = A\sqcup B$ for open $A$ and $B$, then
$B = X\smallsetminus A = A^c$ and $A = X\smallsetminus B = B^c$
are both complements of open sets and thus open.

\begin{example}
	The set $E = (-10, -1) \cup (1, 10)$ with the metric from $\RR$
	is not connected, as $(-10, 1)$ and $(1, 10)$ are clopen in $E$.
\end{example}

\begin{exercise}[Disconnectedness of subspaces]
	\label{wk8:exer:subspace-disconn}
	Given metric space $(X, d)$, show that
	a subset $E\subseteq X$ is disconnected iff
	there exist disjoint sets $A,\,B\subseteq X$ that are open in $X$
	such that $A\cap E\ne\varnothing$, $B\cap E\ne\varnothing$,
	$A\cap B\cap E = \varnothing$, and $E\subseteq A\cup B$.
	Use this to argue that $(-10, -1)\cup (1, 10)\subseteq \RR$
	is disconnected.
\end{exercise}

What should the connected sets in $\RR$ look like?
\begin{exercise}
	Make a guess.
\end{exercise}

It's not hard to guess that, in fact, the connected sets
are just be the intervals.
\begin{theorem}[Connectedness of open intervals]
	\label{wk8:thm:open-connected}
	The open interval $(0, 1)$ is connected.
\end{theorem}
Of course, this generalizes to any $(a, b)$.
\begin{proof}
	Suppose that $(0, 1)$ is a disjoint union of two open sets $A$, $B$.
	Take $x\in A$, and define
	\[U_x = \{\delta > x : [x, \delta)\subseteq A\}\qquad
	V_x = \{\delta < x : (x, \delta]\subseteq A\}.\]
	That is, $U_x$ tells you
	``how far away can you go to the right of $x$
	before reaching a point outside $A$?''
	and $V_x$ encodes the analogous information to the left.
	
	\begin{exercise*}
		Draw a picture!!!
	\end{exercise*}
	
	We can see that
	\[0 \le \inf V_x < x < \sup U_x \le 1\]
	(the middle two inequalities are strict by openness of $A$).
	However, we cannot have both
	$\inf V_x = 0$ and $\sup U_x = 1$, since these two together
	force $A = (0, 1)$.
	\begin{exercise*}
		Prove the above statement.
	\end{exercise*}

	So, suppose WLOG that $\sup U_x < 1$.
	
	\begin{claim}
		$y = \sup U_x$ is not in $A$.
	\end{claim}
	\begin{proof}
		Suppose $y\in A$. Then, we can find an open ball
		centered at $y$, say with radius $\eps$, that is contained in $A$.
		Then, $[x,y)\cup(y,y+\eps) = [x,y+\eps)\subseteq A$,
		so $\sup A > y$, contradiction.
	\end{proof}
	
	This means $y\in B$. But then we can find an open ball
	$(y - \eps, y + \eps)\subseteq B$,
	so $(y-\eps,y+\eps)$ is disjoint from $A$.
	This contradicts the fact that $y = \sup U_x$,
	since this means that $y - \eps$ is an upper bound for $U_x$.
\end{proof}

\begin{remark}
	The idea behind this proof is that
	since open sets in $\RR$ are unions of open intervals,
	we want to somehow look at a ``boundary'' between
	$A$ and $B$, as this is where the issue fundamentally happens
	(just look at how $y$ was an issue regardless of
	whether it was in $A$ or $B$!).
	This is what motivates defining $U_x$ and $V_x$;
	we are essentially looking for a maximal open interval
	$(w, y)$ containing $x$ such that $(w, y)\subseteq A$
	(here $w$ would be $\inf V_x$).
\end{remark}

The similar result for closed intervals is deferred to the homework,
or something.

\begin{theorem}
	\label{wk8:thm:cont-preserve-conn}
	Let $f\colon X\to Y$ be continuous.
	Then, if $E\subseteq X$ is connected, $f(E)\subseteq Y$ is connected.
\end{theorem}

\begin{proof}
	We leverage Exercise \ref{wk8:exer:subspace-disconn}
	to prove the contrapositive.
	Namely, if $f(E)$ is disconnected, then there exist
	open subsets $A$, $B$ of $Y$ that do not intersect inside $f(E)$
	and whose union contains $f(E)$.
	Let $A' = f\inv(A)$ and $B' = f\inv(B)$, which are both open.
	Then, $A'$ and $B'$ are disjoint
	and their union contains $f\inv(f(E))\supseteq E$, which is enough
	to show that $E$ is disconnected.
\end{proof}

\begin{lemma}
	\label{wk8:thm:conn-iff-interval:lemma}
	A set $I\subseteq\RR$ satisfies the condition that,
	for all $a,\,b\in I$,
	\[a\le z\le b\implies z\in I\]
	if and only if it is an interval.
\end{lemma}
\begin{proof}
	If $I$ is an interval then the conclusion is trivial.
	
	For the other direction, let $p = \inf I$ and $q = \sup I$.
	For any $\eps$, we can find $p'$ and $q'$ in $I$ such that
	\[p' < p + \eps,\qquad q - \eps < q',\]
	so $p + \eps \le z \le q-\eps$ implies $z\in I$.
	Taking $\eps$ arbitrarily small gives $p < z < q\implies z\in I$
	(where the inequalities are strict because
	we cannot guarantee that $p$ or $q$ are in $I$).
	This is enough to conclude that $I$ is an interval
	because it is just the set $\{z : p < z < q\}$,
	potentially with $p$ or $q$ added in.
\end{proof}

\begin{theorem}
	\label{wk8:thm:conn-iff-interval}
	$I\subseteq\RR$ is connected if and only if it is an interval.
\end{theorem}
\begin{proof}
	We have shown (in conjunction with the homework, apparently)
	that all intervals are connected.
	
	For the other direction, we prove that
	if $I$ is not an interval then it is disconnected.
	Lemma \ref{wk8:thm:conn-iff-interval:lemma}
	shows that, since $I$ is not an interval,
	we can find $a$ and $b$ in $I$ such that
	there exists a $z$ that is in $[a, b]$ but not $I$.
	Then we can take $A = (-\infty,z)$ and $B = (z,\infty)$
	and use Exercise \ref{wk8:exer:subspace-disconn} to show disconnectedness.
\end{proof}

The upshot?
\begin{theorem}
	Let $(X, d)$ be a connected metric space
	such that $f\colon X\to\RR$ is continuous.
	Then, $f(X)$ is an interval.
\end{theorem}
\begin{proof}
	$f(X)\subseteq\RR$ is connected, so it is an interval
	by Theorem \ref{wk8:thm:conn-iff-interval}.
\end{proof}

\begin{exercise}
	Prove Theorem \ref{wk7:thm:ivt}.
\end{exercise}

\subsection{Path-connectedness}
There is a much more intuitive possible way to define connectedness:

\begin{definition}
	\label{wk8:def:pathconn}
	We say that $a,b\in X$ are \vocab{connected by a path}
	if there exists a continuous function $\gamma\colon[0, 1]\to X$
	such that $\gamma(0) = a$, $\gamma(1) = b$.
	The space $X$ is \vocab{path-connected} if
	any two points in $X$ are connected by a path.
\end{definition}

\begin{theorem}
	\label{wk8:thm:pathconn-implies-conn}
	If $X$ is path-connected, then $X$ is connected.
\end{theorem}
\begin{proof}[Sketch]
	Let $X$ be a union of disjoint open sets $A$, $B$.
	Then, take a path $\gamma$ between $a\in A$ and $b\in B$
	and look at $\gamma\inv(A)$, $\gamma\inv(B)$.	
\end{proof}

Unfortunately the converse of this theorem does not hold.

\pagebreak
\section{Week 9: 10/22 and 10/24}
\subsection{Accumulation point}
\begin{definition}[Accumulation point]\label{wk9:def:accumulationpt}
	Let $E\subseteq X$ be a subset of metric space $X$.
	We say that $x\in X$ is an \vocab{accumulation point} of $E$
	if, for every $\eps > 0$, the ball $B_X(x, \eps)$
	intersects $E$ at a point \alert{besides $x$}.
\end{definition}

In particular $B_X(x,\eps)\cap E$ must be infinite for any $\eps$
in order for $x$ to be an accumulation point.
Of course, accumulation points are limit points,
since we can find a sequence of points converging to $x$
by making $\eps$ smaller and smaller.

The distinction is the following:
consider the set
\[S = \{(x, y)\in\RR^2 : x^2 + y^2 < 1\}\cup\{(0,100)\}.\]
\begin{center}
	\begin{asy}
		draw((-3,0)--(3,0),Arrows(5));
		draw((0,-3)--(0,3),Arrows(5));
		filldraw(circle((0,0),1),red+opacity(0.3),red);
		dot((2.5,0),red);
	\end{asy}
\end{center}

(Diagram not to scale.)
Then, $(100,0)$ is a limit point, but not an accumulation point!!


\subsection{Basic definition of compactness and Weierstrass's theorem}
Last week, we proved the intermediate value theorem
via the machinery of connectedness.
The other part of the two big theorem in week 7 is Weierstrass's theorem.

\begin{definition}[Compactness]\label{wk9:def:compactness}
	Let $(X, d)$ be a metric space.
	We say that $X$ is compact if, whenever $X$ is written as a union
	of open sets
	\[X = \bigcup_{\alpha\in I} U_\alpha,\]
	there exist indices $\alpha_1,\dots,\alpha_n\in I$ such that
	\[X = \bigcup_{i=1}^n U_{\alpha_n}.\]
\end{definition}

\begin{example}[Non-example]
	Note that $X = \RR$ is not compact.
	To prove this, we need to antagonistically choose the $U_\alpha$.
	We can clearly say that
	\[\RR = \bigcup_{\alpha\in[1,+\infty)} U_\alpha, \]
	where $U_\alpha = (-\alpha,\alpha)$.
	Then, if $\alpha_1$, \dots, $\alpha_n$ is a finite set of indices
	such that WLOG $\alpha_1 \le \alpha_2 \le \cdots\le \alpha_n$,
	we would get
	\[U_{\alpha_1}\cup\cdots\cup U_{\alpha_n} = (-\alpha_n,\alpha_n)\ne\RR.\]
	In some sense, the point is that
	we don't have ``too much redundancy''
	with our choice of cover here.
\end{example}

\begin{exercise}\label{wk9:exer:open-not-compact}
	Prove similarly that $(0,1)$ is not compact.
\end{exercise}

Some terminology to make all of this easier to talk about:
\begin{definition}[Covers and subcovers]\label{wk7:def:covers-subcovers}
	We say that a collection of open sets $\mathcal {U} = \{U_\alpha\}$
	is a \vocab{(open) cover} for $E$ if
	\[E\subseteq\bigcup_{U\in\mathcal U} U.\]
	If $\mathcal U$ and $\mathcal U'$ are both covers of $E$
	such that $\mathcal U' \subseteq \mathcal U$,
	then $\mathcal U'$ is a \vocab{subcover} of $\mathcal U$.
	%TODO (do people say ``of $\mathcal U$'' or ``of $E$''?)
	Obviously, we attach the adjective ``finite'' to a (sub)cover
	if the relevant collection of open sets is finite.
\end{definition}

So we can more concisely say
``$X$ is compact if every cover of $X$ has a finite subcover.''
As usual, when talking about subsets $E$ of some ambient metric space $X$,
we work with the induced metric, so $E$ is compact
if every open cover \alert{in $E$} has a finite subcover.
This turns out to not be an issue at all.
\begin{lemma}\label{wk9:lemma:subset-compactness}
	$E$ is compact if and only if
	every open cover of $E$ \alert{in $X$} has a finite subcover.
\end{lemma}

\begin{proof}
	We will just prove the direction assuming $E$ is compact
	(in the induced metric);
	the other direction is left as an exercise. In this case, suppose that
	$\mathcal{V} = \{V_\alpha\}$ is a cover for $E$.
	Then, $U_\alpha = V_\alpha\cap E$ is open in $E$
	by Problem 4 in the homework, so $\mathcal{U} = \{U_\alpha\}$
	is an open cover for $E$ in $E$,
	and then we can extract a finite subcover
	$\{U_{\alpha_i} : 1\le i\le n\}$ whose union is $E$,
	which in turn means we have a finite subcover
	\[\bigcup_{i=1}^n V_{\alpha_i}\supset
	\bigcup_{i=1}^n U_{\alpha_i} = E.\]
\end{proof}

\begin{exercise}
	Prove the other direction of Lemma \ref{wk9:lemma:subset-compactness}.
\end{exercise}

This lemma says something very important:
when talking about the compactness of a set,
we don't need to refer to its parent metric space at all
besides to get our open sets!
This is very different from, say, deciding when sets are open,
which is dependent on whether we are in $X$ or $E$
(for example, $[0,1)$ and $[0,2]$ are open in $[0,2]$ but not $\RR$).

\begin{example}[Compact sets]
	Some examples of compact sets:
	\begin{itemize}
		\ii All finite sets.
		\ii Closed intervals, as we show in
		Theorem \ref{wk9:thm:closedinterval-complete}.
		\ii Any closed and bounded subset of $\RR^n$, in fact ---
		this is called the \hyperref[wk9:thm:heine-borel]{Heine-Borel theorem}.
	\end{itemize}
\end{example}

\begin{theorem}\label{wk9:thm:closedinterval-complete}
	The interval $[0,1]$ is compact.
\end{theorem}

\begin{remark}
	In fact, it turns out that the set of rationals $q$
	for which $0 \le q\le 1$ is not compact!
	So we really have to make use of the fact that
	we are living in $\RR$ here.
\end{remark}

\begin{proof}
	Let $\mathcal{U} = \{U_\alpha\}$ be an open cover of $[0,1]$.
	Now, suppose for the sake of contradiction that we can't
	find a finite subcover here, and define $I_0 = [0,1] = [a_0, b_0]$.
	Then, one of $[0,\half]$ and $[\half,1]$ are also covered by $\mathcal U$.
	One of them must also not
	have a finite subcover, since if both intervals
	had a finite subcover then we could just merge them together.
	Let $I_1 = [a_1, b_1]$ be the interval
	that doesn't have a finite subcover (pick one at random if both are OK).
	
	Then, again, $I_1$ can be split into $[a_1, \frac{a_1+b_1}{2}]$
	and $[\frac{a_1 + b_1}{2}, b_1]$, and yet again
	one of these cannot have a finite subcover,
	which we will call $I_2 = [a_2, b_2]$.
	
	In this way, we inductively define a sequence of intervals $I_n = [a_n,b_n]$
	such that $I_0\supseteq I_1\supseteq I_2\supseteq\cdots$,
	none of the $I_n$ has a finite subcover in $\mathcal U$, and
	$b_n - a_n = 2^{-n}$.
	
	\begin{exercise}
		Prove that $(a_n)$ and $(b_n)$ both converge
		to the same number $x$.
	\end{exercise}
	
	The result of the exercise tells us that
	the intersection of all the $I_n$ is the singleton set
	\[\bigcap_{n\ge0} I_n = [x,x] = \{x\}.\]
	We know that $x$ has to belong to some $U\in\mathcal U$.
	But then this means that $B(x,\eps) = (x-\eps,x+\eps)\subseteq U$
	for some $\eps>0$ because $U$ is open.
	Since $a_n$ and $b_n$ converge to $x$,
	both $a_n$ and $b_n$ must belong to this interval eventually,
	hence $I_n = [a_n,b_n]\subseteq(x-\eps,x+\eps)$ eventually.
	This means that in fact $I_n$ has a finite cover in $\mathcal U$
	eventually, so we have found a contradiction!
\end{proof}

\begin{exercise}
	Where exactly did we use the fact that $\RR$ is complete?
	Use this to argue that $[0,1]\cap\QQ$ is not compact.
\end{exercise}

\begin{theorem}[Preservation of compactness]
	\label{wk9:thm:cont-preserve-comp}
	If $f\colon X\to Y$ is continuous
	and $K\subseteq X$ is compact,
	then $f(K)\subseteq Y$ is also compact.
\end{theorem}

\begin{exercise}
	Prove the theorem above.
	Also, compare to Theorem \ref{wk8:thm:cont-preserve-conn}.
\end{exercise}

\begin{theorem}\label{wk9:thm:compact-is-closed}
	Compact subsets are closed.
\end{theorem}

\begin{proof}
	Let $K$ be compact subset of metric space $X$.
	The goal obviously is to prove that $K$ contains its limit points.
	
	Suppose otherwise, so $\ell\notin K$.
	Then, consider the open sets
	\[U_n = \left\{x\in X : d(x,\ell) > \frac 1n\right\}.\]
	
	\begin{exercise}
		Prove that the $U_n$ really are open in $X$.
	\end{exercise}

	Then the union of the $U_n$ contains every point in $X$ except $\ell$,
	which means the union contains $K$
	because of our assumption that $\ell$ isn't in $K$.
	By compactness, we can find some $n_1$, \dots, $n_k$ such that
	\[K\subseteq U_{n_1}\cup\cdots\cup U_{n_k},\]
	but we can see that this union is just $U_{M}$
	where $M = \max\{n_1, \dots, n_k\}$,
	which is the complement of a closed ball.
	That is, there is a ball $B_X(\ell,\eps)\subseteq X\setdiff K$
	(since the closed ball contains an open ball).
	The contradiction then comes from
	taking a sequence $(x_n)$ of points in $K$ that converges to $\ell$,
	as $(x_n)$ must eventually be in $B_X(\ell,\eps)$.
\end{proof}

You are now free to extract the Weierstrass theorem!
\begin{exercise}
	Prove Theorem \ref{wk7:thm:weierstrass}.
	You will have to show that $[a, b]$
	is the only kind of interval (infinite or not)
	that is compact.
	
	You should also note that the proof that $(0, 1)$ is not compact
	is completely trivial now due to Theorem \ref{wk9:thm:compact-is-closed},
	as opposed to the work
	that you had to do in Exercise \ref{wk9:exer:open-not-compact}.
\end{exercise}

\subsection{More properties of compactness}

\begin{theorem}\label{wk9:thm:prod-compact}
	If $X$ and $Y$ are compact, then so is $X\times Y$.
\end{theorem}

\begin{proof}
	The idea behind this proof is that
	if $U\subseteq X\times Y$ is open, then $U$ is a product of
	sets of the form $B_X(x, \eps)\times B_Y(y,\eps)$.
	\todo{proof}
	
	\begin{center}
		\begin{asy}
			draw((0,0)--(3,0));
			draw((0,0)--(0,3));
			draw((1,3)--(1,0),linewidth(1)+red);
			dot((1,0),red);
			label("$X$", (3,0),dir(0));
			label("$Y$", (0,3),dir(90));
		\end{asy}
	\end{center}
	For some point $x\in X$, we look at the ``line''
	$\{x\}\times Y$.
	Then, we have a continuous map (actually a homeomorphism)
	$f_x\colon Y\to \{x\}\times Y$ given by $y\mapsto (x,y)$,
	which means this line is compact too.
	Fixing a cover $\{U_\alpha\}$, we can say that
	$\{U_{x, i} : 1\le i\le n(x)\}$ is some finite cover of $\{x\}\times Y$.
	Each $U_{x,i}$ can be written as $A_{x,i}\cap B_{x,i}$
	where $A_{x,i}$ and $B_{x,i}$ are open balls in $X$ and $Y$ respectively.
	\todo{finish pf (look at nb)}
\end{proof}

\begin{corollary}\label{wk9:corollary:kcell-compact}
	Given real numbers $a_1$, \dots, $a_n$, $b_1$, \dots, $b_n$,
	the set $[a_1, b_1]\times [a_2, b_2]\times\cdots\times [a_n, b_n]$
	is compact.
\end{corollary}

Another way to generate compact sets:
\begin{lemma}\label{wk9:lemma:closed-subset-is-compact}
	If $X$ is compact and $K\subseteq X$ is closed, then $K$ is compact.
\end{lemma}

\begin{proof}
	Let $\{U_\alpha\}$ be an open cover of $K$.
	Then $\{U_\alpha\}\cup\{X \setdiff K\}$ is an open cover of $X$
	because $X\setdiff K = U_0$ is open in $X$.
	\todo{finish}
\end{proof}

Nicely enough, we can say something much stronger than Corollary
\ref{wk9:corollary:kcell-compact}, fully characterizing
all compact sets in $\RR^n$.
\begin{theorem}[Heine-Borel]\label{wk9:thm:heine-borel}
	A set $K\subseteq\RR^n$ is compact if and only if
	it is closed in bounded.
\end{theorem}

When we say that a set $E$ is \vocab{bounded}, we mean that the quantity
\[\opname{diam}(E) = \sup\{d(x,y):x,y\in E\}\]
is a real number.

\begin{proof}[Proof of Theorem \ref{wk9:thm:heine-borel}]
	It is not hard to prove the forward direction.
	\begin{lemma}
		If $K$ is closed and bounded, it is compact.
	\end{lemma}
	\begin{proof}
		If $K$ is closed and bounded, then $\opname{diam}(K) = R < +\infty$,
	so $K\subseteq B(\mb x,R)$ where $\mb x$ is any point in $K$.
	We know that the standard metric on $\RR^n$ is equivalent to
	$d_{\infty}(x,y) = \max_i|x_i - y_i|$, meaning
	there exists an $R'$ such that
	$B(\mb x, R)\subseteq B_\infty(\mb x, R')$.
	The balls in this new metric are just products of open intervals,
	so, $\ol{B_\infty(\mb x,R')}$ is compact
	and $K$ is a closed subset of this compact set,
	thus it is compact.
	\end{proof}
	
	The other direction uses very intrinsically the idea that
	we are working in a metric space (in fact that we are in $\RR$).
	\begin{definition}\label{wk9:thm:sequential-compactness}
		A metric space is $(X, d)$ is \vocab{sequentially compact} if
		every sequence $(x_n)$ in $X$ has a convergent subsequence
		(the main point is that \alert{the limit is in $X$}).
	\end{definition}

	\begin{lemma}\label{wk9:lemma:seq-compact-if-compact}
		If $K$ is compact, then it is sequentially compact.
	\end{lemma}
	\begin{proof}
		Let $E = (x_n)$ be a sequence of elements in $K$. If $E$ is finite
		then we can find a constant subsequence by pigeonhole,
		so suppose $E$ is infinite.
		
		The goal is to show that $E$ has an accumulation point;
		suppose otherwise for contradiction.
		Then, each $z\in E$ is contained by some open set
		whose intersection with $E$ is just $\{z\}$.
		Take an open cover $U_x$ indexed by points in $X$
		such that, if $z\in E$, then $U_z$ contains no points of $E$
		except $z$ (for example, each $U_z$ could be a sufficiently small
		open ball centered at $z$).
		Then on one hand we can find a finite subcover of $X$ by compactness,
		but on the other hand such a finite subcover
		would have only finitely many elements of the infinite set $E$.
	\end{proof}

	\begin{lemma}
		If $K$ is sequentially compact, then it is closed and bounded.
	\end{lemma}
	\begin{proof}
		Suppose that $K$ is not closed.
		Then, take a limit point $\ell$ of $K$ that is outside of $K$,
		and let $(x_n)$ be a sequence of points in $K$ converging to $\ell$;
		this sequence has no convergent subsequence in $K$
		because each subsequence also converges to $\ell$.
		
		Suppose now that $K$ is unbounded.
		Then, for some arbitrary $x\in K$, consider the sets
		\[S_n = \{y\in K : n-1\le d(x, y) < n\}\]
		for positive integers $n$.
		These sets form a partition of $K$.
		If there existed an $N$ for which $n > N$ implies $S_n$ is empty,
		then $K$ would be bounded, so infinitely many $S_n$ are nonempty,
		say $S_{n_1}$, $S_{n_2}$, \dots. Then take the sequence $(x_{n_k})$
		such that $x_{n_k}\in S_{n_k}$. Each subsequence is unbounded,
		but it is easy to show that a convergent sequence must be bounded
		(exercise!), so $(x_{n_k})$ has no convergent subsequence.
		
		Therefore $K$ must be both closed and bounded
		if it is sequentially compact.
	\end{proof}

	These three lemmas constitute a proof of the theorem!
\end{proof}

\subsection{Uniform continuity}
Recall the definition of continuity of a function $f\colon X\to Y$:
for each $x\in X$, given $\eps > 0$, there exists $\delta > 0$ such that
\[d_X(x, y) < \delta \implies d_Y(f(x), f(y)) < \eps.\]
The thing is that $\delta$ is dependent on both $x$ and $\eps$.
We would like it to be dependent on just $\eps$; that is, we'd like to say
for each $\eps > 0$, there exists a universal $\delta > 0$ such that
\[d_X(x, y) < \eps \implies d_Y(f(x), f(y)) < \eps.\]

\begin{theorem}\label{wk9:thm:compact-domain-impl-unif-cont}
	If the domain of a continuous function $f\colon X\to Y$ is compact,
	then the function is uniformly continuous.
\end{theorem}

\pagebreak
\section{Week 10: 10/29 and 10/31}
We begin by proving Theorem \ref{wk9:thm:compact-domain-impl-unif-cont}.

\begin{proof}[Proof of theorem]
	Suppose $f$ is not uniformly continuous.
	Then, there exists an $\eps > 0$ such that, for each $\delta > 0$,
	there are two points $x,\,y\in X$ such that
	$d_X(x, y) < \delta$ but $d_Y(f(x), f(y)) \ge \eps$.
	
	In other words, for each $n\in\NN$, there exist $x_n$ and $y_n$ such that
	$d_X(x_n, y_n) < \frac 1n$ and $d(f(x_n), f(y_n)) > \eps$.
	Now, since $X$ is compact, it is sequentially compact
	(Lemma \ref{wk9:lemma:seq-compact-if-compact}).
	The trick is that we want convergent subsequences
	$(x_{n_k})$ and $(y_{n_k})$ of $(x_n)$ and $(y_n)$
	\emph{with the same indices}, since then the condition on $x_n$ and $y_n$
	ensures that the limits of these subsequences are equal.
	
	To generate these convergent subsequences, we do a slick trick:
	first, let $(x_{m_k})$ be a convergent subsequence of $(x_n)$.
	Then, take a convergent subsequence of $(y_{m_k})$
	and call it $(y_{n_k})$. The sequence $(n_k)$ is a subsequence of $(m_k)$,
	so $(x_{n_k})$ is also a subsequence of the convergent sequence $(x_{m_k})$,
	thus $(x_{n_k})$ also converges!
	
	Moreover, we can argue (e.g. by squeeze theorem) that
	the limits are the same.
	
	\begin{remark}
		Can you also say that once you know $x_{m_k}\to\alpha$
		you also know $y_{m_k}\to\alpha$ by the distance condition?
	\end{remark}

	But then we have two sequences with the same limit,
	so $f(x_{n_k})$ and $f(y_{n_k})$ are equal,
	which contradicts the fact that their distance is at least $\eps$.
\end{proof}

\subsection{Completeness}
Recall the following definition:
\begin{definition}[Cauchy sequences in $\RR$]
	A sequence of real numbers $(x_n)$ is \vocab{Cauchy} if
	for each $\eps > 0$, there exists an $N$ such that
	\[m,\,n > N\implies |x_m - x_n| < \eps.\]
\end{definition}

Obviously, we can generalize this to arbitrary metric spaces.
\begin{definition}[Cauchy sequences]\label{wk10:def:cauchyseq}
	A sequence of points $(x_n)$ in a metric space $(X, d)$ is \vocab{Cauchy}
	if, for each $\eps > 0$, there exists an $N$ such that
	\[m,\,n > N\implies d(x_m, x_n) < \eps.\]
\end{definition}

\begin{exercise}
	Prove that all convergent sequences are Cauchy.
\end{exercise}

Intuitively, the idea is that a sequence is Cauchy if
``it gets very bunched up as you go further in the sequence,''
i.e. it \emph{should} have a limit (even if it doesn't actually converge).
This motivates the following definition:
\begin{definition}[Completeness]\label{wk10:def:completeness}
	$(X, d)$ is \vocab{complete} if every Cauchy sequence in $X$
	converges in $X$.
\end{definition}

This is one approach to constructing $\RR$, actually:
we can identify each element of $\RR$ with equivalence classes
of Cauchy sequences!

\begin{lemma}\label{wk10:lemma:cauchybounded}
	Cauchy sequences are bounded.
\end{lemma}
\begin{proof}
	The idea is that we can bound cofinitely many of the points of the sequence
	in a ball of radius $\eps$ centered at some $x_n$;
	say all $x_k$ for $k > N$ are captured in this ball.
	Then, it's a matter of growing the ball to have radius bigger than
	\[\max(\eps, d(x_n, x_1), d(x_n, x_2), \dots, d(x_n, x_N)),\]
	at which point we have captured the whole sequence.
\end{proof}

\begin{theorem}\label{wk10:thm:seqcomp-impl-complete}
	If $X$ is compact, then it is complete.
\end{theorem}

\begin{proof}
	We want to show that any Cauchy sequence in $X$ converges in $X$.
	The idea is that if a Cauchy sequence has a convergent subsequence,
	then it is convergent; indeed, letting $\ell$ be a subsequential limit,
	we see that, for any $\eps > 0$,
	\[d(x_n, \ell) \le d(x_{n_k}, x_n) + d(x_{n_k}, \ell) < 2\eps\]
	is true eventually.
	We can now leverage sequential compactness to win instantly.
\end{proof}

\begin{theorem}\label{wk10:thm:Rn-complete}
	$\RR^n$ is complete.
\end{theorem}
\begin{proof}
	Any Cauchy sequence is bounded, so it can be contained in
	a closed open ball, which is compact.
	This lets us see that all Cauchy sequences converge. Yay!
\end{proof}

\begin{theorem}\label{wk10:thm:l2-complete}
	$\ell^2$ is complete.
\end{theorem}

\subsection{Sequences of functions}
\begin{definition}[A first definition]\label{wk10:def:ptwise-conv}
	Let $f_n: I\to\RR$ be a sequence of functions
	and $E$ a subset of $I$.
	We say that $f_n\to f$ \vocab{pointwise} on $E$ if
	$\lim_{n\to\infty} f_n(x) = f(x)$ for all $x\in E$.
\end{definition}

\todo{examples}
\begin{example}
	content
\end{example}

We can rewrite the definition as saying that
for each $\eps > 0$ and $x\in E$,
\[|f_n(x) - f(x)| < \eps\]
is true eventually, say for $n > N$.
Here we run into something similar to the distinction of
continuity and uniform continuity: $N$ depends on both $x$ and $\eps$.
Indeed, we can define \emph{uniform} convergence:
\begin{definition}[Uniform convergence]\label{wk10:def:unif-conv-func}
	$f_n\to f$ \vocab{uniformly} if the threshold $N$
	does not depend on $x$.
\end{definition}

\begin{exercise}
	Show that $f_n(x) = \frac{1}{n^2x^2}$ does not converge to $f\equiv0$
	uniformly on $(0,1)$,
	while $g_n(x) = \frac{\sin(nx^2)}{n}$ does converge uniformly on $(0,1)$.
\end{exercise}

Consider
\[\mathcal B(E,\RR) = \left\{f\colon E\to\RR :
\sup_{x\in E} |f(x)| < +\infty \right\}.\]
Then, we have a metric
\[d(f, g) = \sup_{x\in E} |f(x) - g(x)|.\]
Under this metric, $d(f, g) = \eps$ means that
at each $x\in E$, $f$ and $g$ are at most $\eps$ apart.
This is nice because it encodes uniform convergence as follows:
\begin{lemma}
	$f_n\to f$ uniformly on $E$ iff
	$f_n\to f$ in $\mathcal B(E,\RR)$.
\end{lemma}

\begin{theorem}
	$\mathcal B(E,\RR)$ is complete.
\end{theorem}

\begin{proof}
	Consider a Cauchy sequence of functions $f_n(x)$.
	Then, the sequence of numbers $f_n(c)$ is Cauchy for all $c\in E$,
	so $f_n$ has a pointwise limit $f$.
	This is a bounded function: for $x\in E$,
	we can pick $N$ so that $n > N$ implies $|f_n(x) - f(x)| < 1$
	(by definition of convergence), so, for $n > N$,
	\begin{align*}
		|f(x)| &\le |f_n(x)| + |f_n(x) - f(x)|\\
		&\le \sup_{x\in E} |f_n(x)| + 1 < +\infty.
	\end{align*}

	To show that this converges with respect to our actual metric $d$,
	pick an $\eps > 0$ and take $m$ and $n$ to be integers that are large enough
	that $d(f_n, f_m) < \eps$; then,
	\begin{align*}
		|f_n(x) - f(x)| &\le |f_n(x) - f_m(x)| + |f_m(x) - f(x)|\\
		&\le d(f_n, f_m) + |f_m(x) - f(x)|\\
		\implies |f_n(x) - f(x)| &< \eps + |f_m(x) - f(x)|.
	\end{align*}
	By pointwise convergence we can take the limit as $m\to\infty$
	and take the supremum to get that
	$\sup_{x\in E}|f_n(x) - f(x)| = d(f_n, f) \le\eps$ for all $\eps > 0$.
\end{proof}

\begin{definition}
	We call $C^0(E, \RR)$ the space of continuous functions
	mapping $E\to\RR$,
	and $C^0_b(E,\RR)$ the space of bounded continuous functions
	from $E$ to $\RR$.
\end{definition}

\begin{theorem}\label{wk10:thm:bounded-cont-funcs-closed}
	$C^0_b(E,\RR)$ is a closed subset of $\mathcal B(E,\RR)$.
\end{theorem}

\begin{proof}
	Let $f_n$ be a sequence of continuous bounded functions
	that converge to some $f\in\mathcal B(E,\RR)$.
	We wish to show that $f$ is continuous; let $\eps > 0$.
	By triangle inequality
	\[|f(x) - f(x_0)| \le
	|f(x_0) - f_n(x_0)| + |f_n(x_0) - f_n(x)| + |f_n(x) - f(x)|.\]
	Take $n$ large so that the first and third terms are smaller than $\eps/3$
	(since $f_n\to f$ uniformly --- we bound both by $d(f, f_n)$),
	and then, since $f_n$ continuous, take $|x_0 - x| < \delta$
	so that the middle term is smaller than $\eps/3$.
	\begin{remark}
		The bound of $|f_n(x) - f(x)|$ cannot be done with
		just pointwise convergence because then
		$\delta$ would have to depend on $n$.
	\end{remark}
	Thus $|f(x) - f(x_0)| < \eps$ for $|x - x_0| < \delta$,
	so $f$ is continuous.
\end{proof}

\begin{corollary}
	$C^0_b(E,\RR)$ is complete.
\end{corollary}

\begin{corollary}
	If $E$ is compact, then $C^0(E,\RR)$ is complete.
\end{corollary}
\begin{proof}
	\dots because $C^0(E,\RR) = C^0_b(E,\RR)$ if $E$ is compact.
\end{proof}

The magic of this theory is that we can create entirely new continuous functions
as limits of sequences of pretty simple functions.
The standard example of this is \vocab{power series}:
take real numbers $c_0$, $c_1$, $c_2$, \dots, and define
\[f(x):= c_0 + c_1x + c_2x^2 + c_3x^3 + \cdots =
\lim_{n\to\infty} \sum_{k=0}^n c_kx^k.\]
In other words $f$ is the limit of functions
$f_n(x) = c_0 + c_1x + \cdots + c_nx^n$.
The question, then, is when the limit exists.
For this, the root test wants us to find
\[\alpha = \limsup |c_kx^k|^{1/k} = |x|\cdot\limsup |c_k|^{1/k}
:= |x|\cdot\beta.\]
Root test thus tells us that the limit exists when $|x| < \frac 1\beta$
and doesn't exist when $|x| > \frac1\beta$, so really
$\frac1\beta$ is the thing we care about --- we will call it
$R$, the \vocab{radius of convergence}.
To clarify: $f$ is defined as a pointwise limit when $x\in(-R, R)$,
it might be defined at $x = \pm R$ (since root test is inconclusive),
and it is never defined for $|x| > R$.

\begin{example}[Exponential function]
	The exponential function $\exp$ can be defined as the power series
	\[\exp(x) := \sum_{k=0}^\infty \frac{x^k}{k!}.\]
	It is easy to check that its radius of convergence is all of $\RR$.
	The familiar properties of the exponential function
	are also not too hard to check.
	We define Euler's number $e$ to then be $\exp(1)$.
\end{example}

In view of the theory above, we'd really like this convergence
to be uniform on $(-R, R)$.

\begin{theorem}[Weierstrass $M$-test]
	Let $E\subseteq\RR$. Suppose that $(g_n)$ is a sequence of
	(bounded) functions $E\to\RR$ such that $|g_k(x)| < M_k$ for all $x\in E$.
	Then, if $\sum M_k$ converges, $f_n = \sum_{k=1}^n g_n$ converges uniformly.
\end{theorem}

\begin{proof}[Proof Sketch]
	Prove that $f_n$ is Cauchy by using the fact that
	the sequence of partial sums of $M_k$ is Cauchy.
\end{proof}

\begin{exercise}
	Show that a power series converges uniformly on every compact subset
	of its interval of (pointwise) convergence.
	Show also that we can't improve this to uniform convergence on $(-R, R)$.
\end{exercise}

\subsection{A really cool application}

\begin{theorem}[Space filling curve]
	There exists a function $f\colon[0,1]\to[0,1]^2$
	that is both continuous and surjective.
\end{theorem}

\begin{exercise}
	Show that $f$ cannot \emph{also} be injective ---
	equivalently, prove that $[0,1]$ and $[0,1]^2$ are not homeomorphic.
\end{exercise}

\begin{proof}
	Split the square into nine smaller squares of equal size
	and draw a path through their centers as follows:
	\begin{center}
		\begin{asy}
			for (int i = 0; i < 4; ++i) {
				draw((i,0)--(i,3));
				draw((0,i)--(3,i));
			}
			draw((0.5,0.5)--(0.5,2.5)--(1.5,2.5)--(1.5,0.5)--(2.5,0.5)--(2.5,2.5), red);
			for (int i = 0; i < 3; ++i) {
				for (int j = 0; j < 3; ++j) {
					dot((i+0.5,j+0.5));
				}
			}
		\end{asy}
	\end{center}

	Then, for any partition into $9^k$ squares, we can patch
	these mini-pieces together as follows:
	\missingfigure{text}
	
	Parameterize this path by \todo{param}[placeholder]
	to get maps $f_n\colon[0,1]\to[0,1]^2$.
	
	\begin{claim}
		For all $m$ and $n$,
		\[|f_n(x) - f_m(x)| \le \sqrt2\cdot 3^{-\min(m, n)}.\]
	\end{claim}
	\begin{proof}
		todo\todo{pf of claim}
	\end{proof}

	Then, the limit of the $f_n$ is a surjective continuous function!
\end{proof}

\pagebreak

\section{Week 11: 11/5 and 11/7}

\subsection{Limits of functions}

\begin{definition}[Limits of functions]\label{wk11:def:func-limit-seq}
	Let $E$ be a subset of $X$ and consider $f\colon X\to Y$.
	Also, let $x_0\in X$ be an accumulation point of $E$.
	Then,
	\[\lim_{x\to x_0} f(x) = L\]
	means that for each sequence $(x_n)$ of points in $E$ 
	converging to $x_0$ such that $x_n\ne x_0$ for all $n\in\NN$,
	we have $f(x_n)\to L$.
	To emphasize: the limit does not care about $f(x_0)$.
\end{definition}

The expected limit uniqueness and combination properties hold.

\begin{theorem}\label{wk11:thm:cont-and-real-lims-impl-bounded}
	If $f\colon\RR\to\RR$ is continuous and
	\[\lim_{x\to\infty} f(x) = L_1,\qquad
	\lim_{x\to-\infty} f(x) = L_2\]
	for real $L_1$ and $L_2$, then $f$ is bounded.
\end{theorem}

\begin{proof}
	Suppose for the sake of contradiction that (WLOG)
	\[\sup_{x\in\RR} f(x) = \infty.\]
	Then, we can find a sequence of real numbers $(x_n)$
	such that the sequence $(f(x_n))$ diverges to $\infty$.
	
	\begin{claim}
		$(x_n)$ is unbounded.
	\end{claim}
	\begin{proof}
		If $(x_n)$ were bounded, then we can find a convergent subsequence
		$(x_{n_k})$ with limit $\ell$. By continuity,
		the limit of $(f(x_{n_k}))$ is $f(\ell)\in\RR$,
		but then $(f(x_n))$ has a limit, so the limit must be
		$f(\ell)$, contradicting the fact that the limit is $\infty$.
	\end{proof}
	
	But now we have found a sequence such that $x_n\to\infty$
	or $x_n\to-\infty$ such that $f(x_n)$ blows up to $\infty$,
	which contradicts the fact that $f(x_n)$ should converge to
	either $L_1$ or $L_2$.
\end{proof}

\begin{definition}[Limits of functions, again]\label{wk11:def:func-limit-epsdel}
	Let $(X, d_X)$ and $(Y, d_Y)$ be metric spaces
	and let $f\colon X\to Y$ be a function. Then,
	\[\lim_{x\to x_0} f(x) = \ell\]
	is equivalent to saying that, for every $\eps > 0$,
	there exists a $\delta > 0$ such that
	\[0 < d_X(x, x_0) < \delta \implies d_Y(\ell, f(x)) < \eps.\]
\end{definition}

\subsection{Derivatives}
And now comes the part where we actually have to live in $\RR$.
\begin{definition}[Differentiability]\label{wk11:def:diffable}
	Let $I\subseteq\RR$ be an open interval.
	A function $f\colon I\to\RR$ is \vocab{differentiable}
	at a point $x_0\in I$ if the limit
	\[f'(x_0) := \lim_{x\to x_0} \frac{f(x) - f(x_0)}{x - x_0}\]
	exists in $\RR$.
	We call this limit the \vocab{derivative} of $f$ at $x_0$.
\end{definition}

\begin{exercise}[Equivalent formulation]
	Show that if $f$ is differentiable at $x_0$, then
	\[f'(x_0) = \lim_{h\to 0} \frac{f(x_0 + h) - f(x_0)}{h}.\]
\end{exercise}

\begin{lemma}[Differentiable implies continuous]
	\label{wk11:lemma:diffable-impl-cont}
	If $f\colon I\to\RR$ is differentiable at $x_0$,
	then it is continuous at $x_0$.
\end{lemma}

\begin{proof}
	By limit properties,
	\[\lim_{x\to x_0} f(x) - f(x_0) =
	\lim_{x\to x_0} \frac{f(x) - f(x_0)}{x - x_0}\cdot (x - x_0)
	= f'(x_0)\cdot 0 = 0,\]
	so $f(x)\to f(x_0)$ as $x\to x_0$.
	(Or you can get into the details with sequences
	to prove the same thing.)
\end{proof}

\begin{theorem}[Derivative is zero at extrema]
	\label{wk11:thm:zero-derivative-extrema}
	If $f\colon I\to\RR$ is differentiable at $x_0$ and
	$x_0$ is a local maximum (minimum) of $f$,
	i.e. $f(x) \le f(x_0)$ ($f(x)\ge f(x_0)$) for all $x$ in some
	open ball containing $x$, then $f'(x_0) = 0$.
\end{theorem}

\begin{proof}
	We prove the theorem when $x_0$ is a local maximum,
	i.e. $f(x_0) \ge f(x)$ whenever $|x - x_0| < \eps$.
	Suppose for the sake of contradiction that $f'(x_0) \ne 0$,
	and say WLOG that the derivative is positive.
	Then, there is a $\delta> 0$ such that
	\[\frac{f(x) - f(x_0)}{x - x_0} > 0\]
	whenever $|x - x_0| < \delta$.
	Picking $x$ so that $0 < x - x_0 < \min(\delta, \eps)$,
	this forces $f(x) - f(x_0) > 0 \iff f(x) > f(x_0)$,
	contradiction.
\end{proof}

\begin{theorem}[Rolle's theorem]\label{wk11:thm:rolle}
	Let $f\colon[a,b]\to\RR$ be a continuous function
	that is differentiable on $(a,b)$.
	If $f(a) = f(b)$, then there exists $x\in(a, b)$
	for which $f'(x) = 0$.
\end{theorem}

\begin{proof}
	By Weierstrass's theorem, there exist points $m$ and $M$ such that
	\[f(m) \le f(x) \le f(M)\]
	for all $x\in[a,b]$. If $f$ is constant then the result is trivial
	(the derivative is zero everywhere), so suppose $f$ isn't constant.
	Then, say WLOG $f(M)\ne f(a)$. Since $M$ is a local maximum of $f$,
	Theorem \ref{wk11:thm:zero-derivative-extrema} says that $f'(M)$ is zero.
\end{proof}

\begin{theorem}[Mean value theorem]\label{wk11:thm:mvt}
	If $f\colon[a,b]\to\RR$ is continuous
	and differentiable on $(a,b)$, then there exists an $x\in(a, b)$ such that
	\[\frac{f(b) - f(a)}{b - a} = f'(x).\]
\end{theorem}

\begin{exercise}
	Prove Theorem \ref{wk11:thm:mvt}. (You'll have to use Rolle's theorem.)
\end{exercise}

\begin{corollary}\label{wk11:corollary:zero-derivative-iff-const}
	Prove that if $f\colon (a, b)\to\RR$ is a differentiable function
	then $f'(x) = 0$ for all $x$ iff $f$ is a constant function.
\end{corollary}

\begin{exercise}
	Prove Corollary \ref{wk11:corollary:zero-derivative-iff-const}.
	Prove that if $f,g\colon(a, b)\to\RR$ are differentiable then
	$f'(x) = g'(x)$ for all $x$ iff $f - g$ is a constant function.
\end{exercise}

\begin{corollary}
	Let $f\colon(a,b)\to\RR$ be differentiable.
	\begin{enumerate}[(i)]
		\ii If $f' \ge 0$ then $f$ is increasing.
		\ii If $f' \le 0$ then $f$ is decreasing.
		\ii If the inequalities are strict in (i) and (ii)
		then the monotonicity is also strict.
	\end{enumerate}
\end{corollary}

\begin{proof}
	We just prove that $f' > 0$ implies $f$ strictly increasing.
	This means we want to prove that $y > x$ implies $f(y) > f(x)$.
	Indeed, let $y > x$. There exists a $\xi\in(x, y)$ such that
	\[\frac{f(y) - f(x)}{y - x} = f'(\xi) > 0,\]
	so $y -x > 0$ implies $f(y) - f(x) > 0$ too.
\end{proof}

\subsection{Sequences of functions, again}

\begin{definition}[$C^k$ functions]\label{wk11:def:Ck-funcs}
	A function is said to be \vocab{class $C^k$}
	if it is differentiable $k$ times,
	and all derivatives (in particular the $k$th derivative) are continuous.
	So, for example, $C^0$ functions are just the continuous ones
	and $C^1$ functions have continuous derivative.
\end{definition}

\begin{remark}
	We don't really care about functions with discontinuous derivative.
\end{remark}

Now, as we did last week, we'd like to look at functions as a whole.
We saw as a corollary of \ref{wk10:thm:bounded-cont-funcs-closed}
that $C^0[a,b]$ (the continuous functions $[a,b]\to\RR$) form
a complete metric space.

We know that $C^1[a,b]\subseteq C^0[a,b]$,
and we'd hope that, like $C^0[a,b]$, that this is a complete space.
We are not so lucky.
\begin{itemize}
	\ii There exist sequences $(f_n)$ of $C^1$ functions
	with limit $f$ not in $C^1$.
	\ii There exist sequences $(f_n)$ of $C^1$ function
	 with limit $f$ in $C^1$\dots
	 but $f_n'$ need not converge to $f$.
\end{itemize}

So what do we do??
The idea is to make our metric reflect what we want to control.
\begin{definition}[$C^1$ metric]\label{wk11:def:C1-metric}
	Let $d_{C^0} := d$ be the familiar metric on $C^0$.
	We define a new metric
	\[d_{C^1}(f, g) = d_{C^0}(u, v) + d_{C^0}(u', v').\]
	Of course, we can generalize this to $C^k$ as
	\[d_{C^k}(f, g) = d_{C^0}(u, v) + d_{C^{k-1}}(u', v')
	= \sum_{i=0}^k d_{C^0}(u^{(i)} , v^{(i)}).\]
\end{definition}

\begin{theorem}\label{wk11:thm:new-Ck-complete}
	With this new metric, $C^1[a,b]$ is complete.
\end{theorem}

\begin{proof}
	Given a Cauchy sequence $(u_n)$ in $(C^1[a,b], d_{C^1})$, the point is that
	both $(u_n)$ and $(u_n')$ are Cauchy with respect to $d_{C^0}$.
	If $u_n\to u$ and $u_n'\to v$, we wish to show that
	$u$ is class $C^1$ and $u' = v$.
	
	Fix $x\in[a,b]$. Then, consider the functions
	\[\varphi_n(t) = \frac{u_n(t) - u_n(x)}{t - x},\qquad
	\varphi(t) = \frac{u(t) - u(t)}{t - x}.\]
	We know that
	\begin{align*}
		\lim_{n\to\infty} \varphi_n(t) &= \varphi(t)\\
		\lim_{t\to x} \varphi_n(t) &= u_n'(t)\\
		\lim_{t\to x} \varphi(t) &= u'(t)
	\end{align*}
	pointwise.\todo{wait why does $u'$ exist}
	
	\begin{claim}
		$\varphi_n\to \varphi$ uniformly.
	\end{claim}
	\begin{proof}
		todo\todo{todo}
	\end{proof}
	
	\begin{claim}
		If $\varphi_n\to\varphi$ uniformly on some $E$
		and $x$ is an accumulation point, then
		\[\lim_{n\to\infty} \lim_{t\to x} \varphi_n(t) =
		\lim_{t\to x} \lim_{n\to\infty} \varphi_n(t),\]
		i.e. we can swap the limits
		if the convergence of $\varphi$ is uniform.
	\end{claim}
	\begin{proof}
		Let
		\[\lim_{t\to x }\varphi_n(t) = \ell_n\]
		and suppose that $\ell_n\to\ell$, i.e. the LHS is $\ell$.
		Then, the inner limit on the right is $\varphi(t)$, so
		\[|\ell - \varphi(t)| \le
		|\varphi(t) - \varphi_n(t)| + |\varphi_n(t) - \ell_n| + |\ell_n - \ell|\]
		and each piece can be bounded independently of $\eps$
		as long as $|t - x|$ is small enough.
		\todo{what?}
	\end{proof}
\end{proof}

\begin{exercise}
	Generalize to $C^k[a,b]$.
\end{exercise}

\pagebreak
\setcounter{section}{12}
\section{Week 13: 11/19 and 11/21}

\subsection{Derivatives of power series}
We're going to use the theory from the latter portion of Week 11
to look at derivatives of power series.
The situation is that we have a power series $f$ given by the limit of
\[f_n(x) = a_0 + a_1x + a_2x^2 + \cdots + a_nx^n.\]
on the interval $(-R, R)$ of convergence.
Restricting to the interval $[-R', R']\subseteq(-R, R)$, we look at the sequence
\[f_n'(x) = a_1 + 2a_2x + 3a_3x^2 + \cdots + na_nx^{n-1}.\]
Both $f_n$ and $f_n'$ converge uniformly on $[-R', R']$, so $f_n$
converges in $C^1[-R', R']$.

\begin{exercise}
	Show that the radius of convergence of $f_n'$ is the same as that of $f_n$.
\end{exercise}

Via the theory in Week 11, we can see that the derivative of $f$ is just
the power series we get by differentiating termwise.

You should then be able to easily conclude the following:
\begin{theorem}
	Power series are infinitely differentiable
	on their interval of convergence!
\end{theorem}

Now, we play!
\begin{example}[Geometric series]
	We know that $f(x) = \sum_{n=0}^\infty x^n = \frac{1}{1-x}$. Then,
	\[f'(x) = \frac{1}{(1-x)^2} =
	\sum_{n=1}^\infty nx^{n-1},\]
	and you can spam derivatives to get more and more facts like this.
\end{example}

\begin{example}[Taylor series]
	Let $f(x) = a_0 + a_1x + a_2x^2 + \cdots$ be a power series.
	Then, we have
	\[f'(x) = 1\cdot a_1 + 2a_2x + 3a_3x^2 + \cdots,\]
	which, evaluated at $0$, gives $1\cdot a_1$.
	Similarly
	\[f''(0) = 2\cdot 1\cdot a_2,\quad f^{(3)}(0) = 3\cdot2\cdot1\cdot a_3,\]
	and in general $f^{(n)}(0) = n!\cdot a_n$.
	Thus, we can write
	\[a_n = \frac{f^{(n)}(0)}{n!}.\]
	In particular,
	\[f(x) = \sum_{n=0}^{\infty} \frac{f^{(n)}(0)}{n!}\cdot x^n.\]
\end{example}

The caveat is that
\[\limsup\abs{\frac{f^{(n)}(0)}{n!}}^{1/n}\]
need not behave nicely.

\begin{exercise}
	Find an example where the Taylor series
	has a radius of convergence of $0$.
\end{exercise}

\begin{lemma}
	Let $f\colon(a, b)\to\RR$ and let $c\in(a, b)$.
	Suppose also that $f\in C^n(a, b)$. Define
	\[P_k(x) = \sum_{j=0}^k \frac{f^{(j)}(c)}{j!}\cdot (x-c)^j,\]
	and let
	\[R_n(x) = f(x) - P_{n-1}(x).\]
	Then, for a fixed $x$, there exists a $y$ between $c$ and $x$ such that
	\[R_n(x) = \frac{f^{(n)}(y)}{n!}\cdot (x-c)^n.\]
\end{lemma}

\begin{proof}
	This should be screaming MVT; we need to find
	a suitable function to use MVT on.
	We can write
	\[f(x) - P_{n-1}(x) = \frac{M\cdot (x-c)^n}{n!}\]
	for some constant $M$. Then, consider
	\[g(t) = P_{n-1}(t) + \frac{M\cdot(t-c)^n}{n!} - f(t).\]
	In fact, all derivatives (up to the $n$th) of $g$ are
	zero at $c$ and some point between $c$ and $x$,
	so we can look at the $n$th derivative and win \todo{write up thoroughly}
\end{proof}

\subsection{Integration!!!!}

\begin{warning*}
	I was half asleep on the days that the integration lessons were taught,
	so these notes are a post-lesson rendition based primarily on
	my own knowledge. As such, some of the notation and content here
	may not align with the stuff shown in class.
\end{warning*}

As we all know from high school, the goal of integration is to
make the idea of ``area under a graph'' more precise.

In high school geometry, we learned how to compute areas of simple shapes,
like rectangles and triangles; to compute areas of more complicated shapes,
we try to decompose them into these simpler shapes.

\missingfigure{probably can make a neat diagram}

However, this is hard (impossible) for fucky shapes like circles,
so we resort to approximations from inside and outside.

\missingfigure{approximating circles w ngons}

Let's start with extra basic functions.
The most basic function we can come up with is
\[\varphi_{(a,b)}(x) = \begin{cases}
	1 & a < x < b\\
	0 & \text{otherwise}
\end{cases}\]
i.e. the indicator function of the interval $(a, b)$.
More generally, for a (finite) interval $I$, we write
\[\varphi_I(x) = \begin{cases}
	1 & x\in I\\
	0 & \text{otherwise}
\end{cases}\]
We call these \vocab{characteristic} or \vocab{indicator} functions.
(In class, we also use \vocab{simple} functions,
but I don't like that terminology.)

\begin{center}
	\begin{asy}
		draw((-5,0)--(5,0));
		draw((-2,1)--(3,1), red+linewidth(1));
		draw((-5,0)--(-2,0), red+linewidth(1));
		draw((3,0)--(5,0), red+linewidth(1));
		fill((-2,0)--(3,0)--(3,1)--(-2,1)--cycle,red+opacity(0.2));
	\end{asy}
\end{center}

The area under the graph of $\varphi_{(a,b)}$
should obviously be $b-a$, so we write
\[\int_{-\infty}^\infty \varphi := \ell((a,b)) = b-a.\]

\begin{definition}[Step function]
	Consider some finite set of intervals $I_1$, $I_2$, \dots, $I_n$
	and some real constants $c_1$, $c_2$, \dots, $c_n$.
	We call the function
	\[\varphi := \sum_{k=1}^n c_k \varphi_{I_k},\]
	a \vocab{step function}.
\end{definition}

\begin{exercise}
	Show that step functions form a vector space.
	That is, show that (finite) linear combinations of step functions
	are also step functions.
\end{exercise}

We would like to define the integral of $\varphi$ in the obvious way, namely
\[\int_{-\infty}^\infty \varphi =
\sum_{k=1}^n c_k\int_{-\infty}^\infty \varphi_{I_k}.\]
In order for this definition to be sound, we have to check the following:
\begin{exercise}
	Let $\varphi$ be a step function,
	with value $c_k$ on the interval $I_k$ for $1\le k\le n$.
	Then, the expression
	\[\sum_{k=1}^n c_k\int_{-\infty}^\infty \varphi_{I_k}\]
	is invariant of the choice of $I_k$s.
	
	(Hint: given two sets of intervals $I_1$, \dots, $I_n$
	and $J_1$, \dots, $J_m$, consider a ``common refinement''
	of the intervals such that each $I_a$ and $J_b$ can be written as
	a disjoint union of the intervals in the refinement.)
\end{exercise}

This behaves as we would expect the integral to behave.
Specifically:
\begin{exercise}[Step function integral properties]
	Show that the integral is a linear function on the set of step functions.
	That is,
	\[\int_{-\infty}^\infty c_1\varphi_1 + c_2\varphi_2 =
	c_1\int_{-\infty}^\infty \varphi_1 + c_2\int_{-\infty}^\infty \varphi_2\]
	for any step functions $\varphi_1$, $\varphi_2$ and
	real numbers $c_1$, $c_2$.
	Show also that if $\varphi_1\ge \varphi_2$ on all of $\RR$
	then their integrals obey the same inequality.
\end{exercise}

We now have a large enough set of functions that
we can start doing approximation.

\begin{definition}[Darboux integral]
	Let $f\colon\RR\to\RR$ be a bounded function with compact support
	(i.e. the subset of $\RR$ on which $f(x)\ne 0$
	can be contained in some compact set).
	Then, let $S_{top}$ be the set of step functions $\varphi$
	for which $\varphi \ge f$ on all of $\RR$
	and $S_{bot}$ be the set of step functions $\varphi$
	for which $\varphi \le f$ on all of $\RR$.
	We define the \vocab{upper Darboux integral} of $f$ to be
	\[\ol{\int_{-\infty}^\infty} f :=
	\inf\left\{\int_{-\infty}^\infty \varphi : \varphi\in S_{top}\right\}\]
	and the \vocab{lower Darboux integral} of $f$ to be
	\[\underline{\int_{-\infty}^\infty} f :=
	\sup\left\{\int_{-\infty}^\infty \varphi : \varphi\in S_{bot}\right\}.\]
	If these two values coincide, then we can define the
	\vocab{Darboux integral} of $f$ as the common value
	of its upper and lower Darboux integrals.
\end{definition}

\begin{remark}
	In your calculus class, you probably saw Riemann integration,
	where you integrate over an interval $[a, b]$ by picking
	some partition of the interval $a = x_0 \le x_1\le\cdots\le x_n = b$
	and then picking arbitrary values $x_{k-1}\le x_k^*\le x_k$
	so that you can estimate your function on $[x_{k-1}, x_k]$
	with the constant value $f(x_k^*)$.
	Then, we can compute the integral of this approximation
	and take the limit as the subintervals of the partition
	become arbitrarily small.
	This turns out to be equivalent to integration in the sense of Darboux,
	so we need not worry about using ``Darboux integrable''
	and ``Riemann integrable'' interchangeably.
\end{remark}

\begin{exercise}
	Prove that the integral properties from Exercise 13.10
	extend to the Darboux integral.
\end{exercise}

\begin{exercise}
	Show that unbounded functions are not Darboux integrable.
\end{exercise}

\pagebreak
The following lemma makes it much easier to show
whether a function is integrable.
\begin{lemma}[Integrability criterion]\label{wk13:integrability-criterion}
	Let $f\colon\RR\to\RR$ be a bounded function with compact support.
	Then, $f$ is integrable if and only if, for each $\eps > 0$,
	there exist step functions
	$\varphi_1 \le f \le \varphi_2$ such that
	\[\int_{-\infty}^\infty \varphi_2 - \varphi_1 \le \eps.\]
\end{lemma}

\begin{proof}
	If $f$ is integrable, then there is a step function $\varphi_1$
	bounded above by $f$ that is at most $\eps/2$ away from the integral of $f$
	and a step function $\varphi_2$ bounded below by $f$
	that is at most $\eps/2$ away from the integral of $f$,
	by definition of sup/inf;
	thus the integral of the difference between $\varphi_1$ and $\varphi_2$
	is at most $\eps$.
	
	Conversely, if we can always find step functions like described,
	then, for each $n$, then the idea is to construct
	an increasing sequence of step functions below $f$
	and a decreasing sequence of step functions above $f$
	and then conclude that the limit of the increasing sequence
	and decreasing sequence of integrals are the same.\todo{details}
\end{proof}

Now, we can finally show that a fairly large class of functions is integrable.
First, let's define integration on an interval.
\begin{definition}
	The integral of a bounded function $f\colon\RR\to\RR$
	on an interval $[a,b]$ is equal to
	\[\int_a^b f := \int_{-\infty}^\infty f\cdot \varphi_{[a,b]},\]
	recalling that $\varphi_I$ is the indicator function of the interval $I$.
\end{definition}

\begin{theorem}[Increasing functions are integrable]
	Let $f\colon\RR\to\RR$ be increasing. Then $f$
	is integrable on any interval $[a,b]$.
\end{theorem}

\begin{theorem}[Continuous functions are integrable]
	Let $f\colon\RR\to\RR$ be continuous. Then $f$
	is integrable on any interval $[a,b]$.
\end{theorem}

\begin{proof}
	Pick $\eps > 0$. Then, since $[a,b]$ is compact,
	$f$ is uniformly continuous on $[a,b]$.
	Thus there exists a $\delta> 0$ such that
	we can find $a = x_0 < x_1 < \cdots < x_n = b$
	satisfying $x_i - x_{i-1} < \delta$ for $1\le i\le n$ and
	the minimum $m_i$ of $f$ over the interval $[x_{i-1}, x_i]$
	differs from the maximum $M_i$ by at most $\frac{\eps}{b-a}$.
	If we take the step function $\phi$ that is $m_i$ on $[x_{i-1}, x_i]$
	and the step function $\psi$ that is $M_i$ on $[x_{i-1}, x_i]$, then
	$\phi \le f \le \psi$ on $[a,b]$ and
	\[\int_a^b \psi - \phi = \sum_{i=1}^n (M_i - m_i)(x_i - x_{i-1})
	\le \eps\sum_{i=1}^n x_i - x_{i-1} = \frac{\eps}{b-a}\cdot(b-a) = \eps.\]
	We are then done by Lemma \ref{wk13:integrability-criterion}.	
\end{proof}

We can also finally prove the fundamental theorem of calculus.
\begin{lemma}[Mean value theorem for integrals]\label{wk13:integralmvt}
	Let $f\colon[a,b]\to\RR$ be a continuous function.
	Then, there is some $c\in[a,b]$ for which
	\[f(c) = \frac{1}{b-a}\int_a^b f(x)\dif x.\]
\end{lemma}

\begin{proof}
	Bound the integral below by $m := \min_{x\in[a,b]} f(x)$
	and above by $M := \max_{x\in[a,b]} f(x)$.
	The existence of $c$ is then guaranteed by the intermediate value theorem.
\end{proof}

\begin{theorem}[First fundamental theorem of calculus]
	Let $f\colon\RR\to\RR$ be a continuous function.
	For a fixed $a\in\RR$, define
	\[F(x) := \int_a^x f(t)\dif t.\]
	Then, $F$ is differentiable with $F'(x) = f(x)$ for all $x\in\RR$.
\end{theorem}

\begin{proof}[Proof Sketch]
	Fix some $x\in\RR$. Then, if $h > 0$,
	\[\frac{F(x+h) - F(x)}{h} = \frac 1h\int_x^{x+h} f(t)\dif t.\]
	By Lemma \ref{wk13:integralmvt}, the right-hand side is equal to
	$f(\alpha_h)$ for some $\alpha_h\in[x, x+h]$.
	As $h\to 0^+$, the $\alpha_h$ approach $x$, so $f(\alpha_h)$ goes to
	$f(\alpha)$. We can get an analogous result for when $h < 0$,
	so the limit of the difference quotient ends up just being $f(x)$.
\end{proof}

\setcounter{section}{14}
\section{Week 15}
\begin{theorem}[Second fundamental theorem of caluclus]
	If $G\colon[a,b]\to\RR$ is such that $G'(x) = f(x)$
	for all $x\in[a,b]$ and $f$ is Riemann integrable on $[a,b]$, then
	\[\int_a^b f(t)\,dt = G(b) - G(a)\]
\end{theorem}

\begin{proof}
	Consider $F(x) = \int_a^x f(t)\,dt$ and set
	\[H(x) := F(x) - G(x).\]
	\begin{claim}
		$H$ is constant on $[a,b]$.
	\end{claim}
	\begin{proof}
		We can compute
		\[H'(x) = F'(x) - G'(x) = f(x) - f(x) = 0\]
		for all $x\in[a,b]$.
		\todo{wait why is $F$ dif'able? we don't know that $f$ is cont}
	\end{proof}
	
	Then, we get
	\[F(a) - G(a) = H(a) = H(b) = F(b) - G(b),\]
	or equivalently $F(b) - F(a) = G(b) - G(a)$, and then
	\[G(b) - G(a) = F(b) - F(a)
	= \int_a^b f(t)\,dt - \int_a^a f(t)\,dt = \int_a^b f(t)\,dt\]
	as desired.
\end{proof}

Recall that $\mathcal B([a,b],\RR)$ is the space of bounded functions
$[a,b]\to\RR$, which is a complete metric space under the standard sup metric.

Consider the subset $X\subseteq \mathcal B([a,b],\RR)$
consisting of the space of integrable functions.

\begin{theorem}
	$X$ is closed (hence complete).
\end{theorem}
\begin{proof}
	Let $f_n$ be a sequence of functions in $X$ converging uniformly to some $f$.
	We can bound
	\[|f - f_n| \le \frac{\eps}{2(b-a)}\]
	for large enough $n$, independently of $x$. Then,
	upper and lower Darboux integrals are monotonic, so
	\[\underline{\int_a^b} f_n(x) - \frac{\eps}{2(b-a)}\dif x
	\le \underline{\int_a^b} f(x)\dif x
	\le \overline{\int_a^b} f(x)\dif x
	\le \overline{\int_a^b} f_n(x) + \frac{\eps}{2(b-a)}\dif x.\]
	Thus the difference between the upper and lower integrals
	of $f$ over $[a, b]$ is bounded by $\eps$
	(since the upper and lower integrals of $f_n$ are equal by hypothesis),
	i.e. the upper and lower integrals coincide, so
	uniform limits of Riemann integrable functions are still Riemann integrable.
\end{proof}

\begin{corollary}
	Let $f$ be a power series with interval of convergence $(-R, R)$.
	Then $f$ is integrable on every compact interval contained in $(-R, R)$.
\end{corollary}

\begin{proof}
	Power series converge uniformly on compact subsets
	of their interval of convergence, so we can appeal to Theorem 15.2.
\end{proof}

Consider now the map $T\colon X\to\RR$ given by
$T(u) = \int_a^b f(t)\dif t$.

\begin{theorem}
	$T$ is a Lipschitz function (hence uniformly continuous).
\end{theorem}

\begin{proof}[Sketch]
	Just write out the definition of Lipschitz function;
	use linearity and triangle inequality to bound
	\[|T(u) - T(v)| \le (b-a)d(u, v).\qedhere\]
\end{proof}

For integrals of functions that are bounded but not compactly supported, we define
\[\int_0^\infty f(x)\dif x := \lim_{a\to\infty} \int_0^a f(x)\dif x\]
and similarly for integrals from $-\infty$ to $0$, and then we glue them together
for something like
\[\int_{-\infty}^\infty f(x)\dif x :=
\lim_{T\to\infty} \int_0^T f(x)\dif x + \int_{-T}^0 f(x)\dif x.\]

\subsection{Differential equations}
Let $\Omega\subseteq\RR^2$ be open
(for simplicity, let $\Omega = I\times\RR$ for some open interval $I$)
and consider a function $g\colon\Omega\to\RR$ that is continuous.
We would like to look for functions $u\colon\RR\to\RR$ such that
\[u'(t) = g(t, u(t))\]
with a prescribed value $u(t_0) := u_0$.
For example, we can have $g(t, u) := \sqrt{t^2 + u}$ so that we want
\[u'(t) = \sqrt{t^2 + u(t)}\]
and we can prescribe $u(0) = 1$.

In many cases, we at least would like to know that a solution $u$ exists,
even if it isn't actually tractable.

The game plan is as follows. Suppose for a moment that we have a solution $u$.
Then, we can integrate so get
\[\int_0^x g(t, u(t))\dif t = \int_{t_0}^x u'(t)\dif t  = u(x) - u(t_0)
= u(x) - u_0,\]
so we get the rearrangement
\[u(x) = u_0 + \int_{t_0}^x g(t, u(t))\dif t.\]
We know a priori that $u$ is differentible, hence continuous,
so the integral in this expression above is a class $C^1$ function.
We can thus deduce that if a function solves this equation
then it must be class $C^1$.
(In general if $g$ is class $C^k$ then we can deduce that
$u$ is class $C^{k+1}$).

\begin{theorem}
	There exists a $\delta > 0$ and a $C^1$ function
	$u\colon(-\delta + t_0, \delta + t_0)\to\RR$ that satisfy
	the constraints we want (namely $u'(t) = g(t, u(t))$ and $u(t_0) = u_0$),
	if $g$ is $C^1$.
\end{theorem}

\begin{proof}
	Let $X_\delta$ be the set of continuous bounded functions
	$u\colon(-\delta + t_0, \delta + t_0)\to\RR$,
	which we know is complete.
	Consider the operator $T\colon X_\delta\to X_\delta$ defined as
	\[u\mapsto Tu(x) := u_0 + \int_{t_0}^x g(t, u(t))\dif t.\]
	
	We will show that
	\begin{center}
		\framebox{The operator $T$ has a fixed point.}
	\end{center}
	
	This will be done via the following theorem.
	\begin{theorem*}
		[Banach-Cacciopolli fixed point theorem]
		Let $(X, d)$ be a complete metric space and suppose that
		$T\colon X\to X$ is a \vocab{contraction},
		meaning that it is Lipschitz and that its
		Lipschitz constant $L$ is less than $1$.
		Then, $T$ has a \emph{unique} fixed point.
	\end{theorem*}
	
	\begin{proof}
		Let $u_0$ be a point in $X$
		and consider the sequence $(u_n) = (u, Tu, T^2u, \dots)$ given by
		$u_0 := u$ and $u_{n+1} = Tu_n$ for $n\ge 0$. Then,
		\begin{align*}
			d(u_n, u_{n+1}) &= d(Tu_{n-1}, Tu_n)\\
			&\le Ld(u_{n-1}, u_n)\\
			&\le L^2d(u_{n-2}, u_{n-1})\\
			&\vdots\\
			&\le L^nd(u_0, u_1).
		\end{align*}
		Defining $C := d(u_0, u_1)$, we have that
		$d(u_n, u_{n+1}) \le CL^n$ for $L < 1$.
		This lets us estimate $d(u_n, u_m)$ from above by
		$d(u_n, u_{n+1}) + d(u_{n+1}, u_{n+1}) + \dots$ so that
		\[d(u_n, u_m) \le C\cdot (L^n + L^{n+1} + \cdots)
		= \frac{CL^n}{1 - L},\]
		and we get the same bound with $m$, so as long as $\min(m, n)$
		is large enough, the distance between the terms
		is going to be smaller than any given $\eps > 0$.
		Thus the $u_n$ are Cauchy, so they converge (by completeness).
		
		This means
		\[\lim_{n\to\infty} u_n = \lim_{n\to\infty} T(u_{n-1})
		= T(\lim_{n\to\infty} u_{n-1}) = T(\lim_{n\to\infty} u_n),\]
		so the limit of the $u_n$ is a fixed point.
		To show uniqueness, we can see that if $u$ and $v$
		are two distinct fixed points of $T$, then
		\[d(u, v) = d(Tu, Tv) \le L\cdot d(u, v) < d(u, v),\]
		which is a contradiction.
	\end{proof}
	
	To finish, we claim that
	\begin{claim}
		The operator $T\colon X_\delta\to X_\delta$ is Lipschitz.
	\end{claim}
	\begin{proof}
		For two functions $u$, $v$,
		\[|Tu(x) - Tv(x)| = \abs{\int_{t_0}^x g(t, u(t)) - g(t, v_t)\dif t}
		\le \int_{t_0}^x |g(t, u(t)) - g(t, v(t))|\dif t.\]
		\todo{proof}
	\end{proof}
\end{proof}




















\end{document}